% Responsable: Fernando Hernández
\section{La infografía final y conclusiones}

\subsection{Resultado}

La figura~\ref{fig:infografia} es la infografía tal como se entrega. Se genera
completa con \texttt{src/infografia.py} a partir de las tablas de
\texttt{outputs/}, así que cada cifra de la página se puede rastrear hasta el
script que la calcula y, desde ahí, hasta las series del Laboratorio 1. Se exporta
en dos versiones: un PNG de 1650 $\times$ 2550 píxeles (150 ppp), por encima del
mínimo de 1200 $\times$ 2400, y un PDF vectorial de exactamente una página de
11 $\times$ 17 pulgadas para impresión.

\begin{figure}[p]
\centering
\includegraphics[height=0.9\textheight]{infografia.png}
\caption{Infografía final.}
\label{fig:infografia}
\end{figure}

\subsection{Del boceto a la página}

La página respeta el boceto de la sección~\ref{sec:recorrido}: las cinco franjas,
sus proporciones y el recorrido en Z por los bloques 3 y 4. Al construirla
cambiaron cuatro cosas, y cada cambio salió de revisar los datos con más cuidado:

\begin{itemize}[leftmargin=*]
  \item \textbf{La anotación de la pandemia.} El borrador decía ``cero
  visitantes'' en abril--agosto de 2020. Al graficarlo vimos que el total no llega
  a cero, baja a unos 10 mil al mes; solo América del Norte y Europa marcan cero.
  La etiqueta final dice $-$97\% frente a los mismos meses de 2019.
  \item \textbf{El alcance del cambio de metodología.} En el avance creíamos que el
  quiebre de 2023 solo afectaba al total y a América del Centro. Separar los
  visitantes por residencia mostró que también afecta a las tres fronteras. Por
  eso el crecimiento del cuadrante y la primera recomendación usan solo
  extranjeros.
  \item \textbf{El orden de las filas del cuadrante.} Pasó de ordenar cada columna
  de mayor a menor a un orden fijo de series, para que cada serie se pueda seguir
  a lo largo de las cuatro preguntas.
  \item \textbf{La respuesta de estacionalidad.} Se marcan dos series en acento,
  América del Norte y Europa, porque la diferencia entre ellas es de 0.002.
\end{itemize}

\subsection{Cómo responde la página a lo que pide el enunciado}

\begin{center}
\begin{tabularx}{\textwidth}{@{}lX@{}}
\toprule
Requisito & Dónde se cumple \\
\midrule
Titular periodístico y subtítulo & Bloque 1: ``Guatemala recupera a sus viajeros: los
mercados lejanos ya superan 2019''. \\
3 a 4 tarjetas KPI & Bloque 1: cuatro cifras de 32 pt con su línea de contexto. \\
Línea temporal con eventos & Bloque 2: récord de 2019, cierre de 2020, reapertura,
cambio de metodología de 2023 y semestre parcial de 2026. \\
Estacionalidad y fechas que marcan la pauta & Bloque 3, izquierda: mapa de calor con
Semana Santa, vacaciones y bono 14, fin de año y aguinaldo. \\
Cuadrante comparativo & Bloque 3, derecha: las cuatro preguntas con la serie que
responde cada una en acento. \\
Comparación de modelos & Bloque 4, izquierda: MAE de ocho variantes, con el RMSE del
ganador en la nota. \\
Pronóstico con intervalo rotulado ``Estimación'' & Bloque 4, derecha: bandas del 80\% y
95\% sobre la prueba y doce meses proyectados. \\
3 a 5 recomendaciones & Bloque 5: cuatro acciones, cada una con su cifra de respaldo. \\
Ficha técnica y créditos & Pie: fuente INGUAT/IGM, periodo, definiciones, partición de
los modelos y autores. \\
\bottomrule
\end{tabularx}
\end{center}

\subsection{Verificaciones para que los gráficos no mientan}

\begin{itemize}[leftmargin=*]
  \item Todos los ejes de cantidades empiezan en cero: la línea de tiempo, las
  barras de modelos, las del cuadrante y el pronóstico.
  \item Cada columna del cuadrante tiene su propia escala y su unidad escrita en
  el encabezado, para que nadie compare barras de preguntas distintas.
  \item Las cifras afectadas por el cambio de metodología se calculan sin ese
  quiebre (solo extranjeros) o se marcan en la página.
  \item Lo que no es dato observado, la estimación y el semestre parcial, se
  distingue con línea punteada y rótulo propio.
  \item Las fechas del calendario se presentan como coincidencias con los picos y
  no como causas, porque el laboratorio midió el patrón y no su causa.
\end{itemize}

\subsection{Conclusiones}

La pregunta del enunciado tiene dos mitades, y la infografía las responde así.

\paragraph{Qué nos enseñan 17 años de datos.} El turismo hacia Guatemala creció de
forma sostenida entre 2009 y 2019, de 1.88 a 4.13 millones de visitantes al año, y
lo hizo con un calendario muy estable: diciembre, Semana Santa y julio se repiten
cada año. La pandemia lo frenó por completo en 2020, y las fronteras terrestres
con El Salvador fueron las que más tardaron en volver. La lectura más importante
es la que no se ve en el total crudo: contando solo extranjeros, en 2025 llegaron
138 visitantes por cada 100 de 2019. El turismo no está rezagado; la base dejó de
contar a una parte de los viajeros.

\paragraph{Qué debemos anticipar.} El modelo ganador estima unos 3.3 millones de
visitantes entre julio de 2026 y junio de 2027, con picos en julio y diciembre.
Pero el resultado más útil del bloque 4 no es esa cifra, sino que un pronóstico
solo sirve si se actualiza cada mes: así se equivoca tres veces menos.
San Cristóbal es la serie que más crece y la menos predecible, así que es donde
el pronóstico y la capacidad operativa necesitan más margen.

\paragraph{Sobre el diseño.} Cada decisión de la página sale de un principio de
la teoría vista en clase o de una revisión de los datos. Por eso el documento y la
infografía cuentan la misma historia: el contexto lleva a la evidencia, la
evidencia al pronóstico y el pronóstico a la acción.
